
\documentclass{article}
\usepackage{colortbl}
\usepackage{makecell}
\usepackage{multirow}
\usepackage{supertabular}

\begin{document}

\newcounter{utterance}

\centering \large Interaction Transcript for game `cladder', experiment `full\_v1.5\_default', episode 4339 with selfrefine{-}llama{-}3.1{-}8b{-}instruct.
\vspace{24pt}

{ \footnotesize  \setcounter{utterance}{1}
\setlength{\tabcolsep}{0pt}
\begin{supertabular}{c@{$\;$}|p{.15\linewidth}@{}p{.15\linewidth}p{.15\linewidth}p{.15\linewidth}p{.15\linewidth}p{.15\linewidth}}
    \# & \multicolumn{2}{c}{Player} && \multicolumn{2}{c}{Game Master} \\
    \hline

    \theutterance \stepcounter{utterance}  
    & & & \multicolumn{4}{p{0.6\linewidth}}{
        \cellcolor[rgb]{0.9,0.9,0.9}{
            \makecell[{{p{\linewidth}}}]{
                \texttt{\tiny{[P1$\langle$GM]}}
                \texttt{You are an expert in causal inference. The following question is not a typical commonsense query, but rather a meticulously designed question created by a professor specializing in causal inference, intended to assess the students' mastery of the course content.} \\
\texttt{Imagine a self{-}contained, hypothetical world with only the following conditions, and without any unmentioned factors or causal relationships: Maternal smoking status has a direct effect on infant's birth weight and freckles. Health condition has a direct effect on infant's birth weight and freckles. Infant's birth weight has a direct effect on freckles. Health condition is unobserved. Method 1: We look at how maternal smoking status correlates with freckles case by case according to infant's birth weight. Method 2: We look directly at how maternal smoking status correlates with freckles in general. To understand how maternal smoking status affects freckles, is it more correct to use the Method 1 than Method 2?} \\
\texttt{Start your answer with 'Yes' or 'No', followed by additional reasoning or evidence to support your explanation.} \\
            }
        }
    }
     \\ \\

    \theutterance \stepcounter{utterance}  
    & \multicolumn{4}{p{0.6\linewidth}}{
        \cellcolor[rgb]{0.561,0.737,0.561}{
            \makecell[{{p{\linewidth}}}]{
                \texttt{\tiny{[P1$\rangle$GM]}}
                \texttt{No.} \\
\\ 
\texttt{In this hypothetical world, we have a complex causal structure involving direct effects between maternal smoking status, infant's birth weight, and freckles. The presence of unobserved health condition further complicates the analysis.} \\
\\ 
\texttt{Method 1, which conditions on infant's birth weight, is an attempt to control for the confounding effect of birth weight on the relationship between maternal smoking status and freckles. However, this approach is problematic because birth weight is a mediator in the causal pathway from maternal smoking status to freckles. Conditioning on a mediator can lead to biased estimates of the causal effect of interest (Pearl, 2009).} \\
\\ 
\texttt{In contrast, Method 2 directly examines the correlation between maternal smoking status and freckles without conditioning on any intermediate variables. This approach is more likely to yield unbiased estimates of the causal effect of maternal smoking status on freckles, as it does not involve conditioning on a mediator.} \\
\\ 
\texttt{To see why, consider the causal diagram:} \\
\\ 
\texttt{Maternal Smoking Status → Infant's Birth Weight → Freckles} \\
\texttt{Maternal Smoking Status → Freckles} \\
\texttt{Health Condition → Infant's Birth Weight → Freckles} \\
\texttt{Health Condition → Freckles} \\
\\ 
\texttt{When we condition on infant's birth weight (Method 1), we are effectively removing the direct effect of maternal smoking status on freckles, which is a part of the causal effect of interest. This can lead to biased estimates of the causal effect.} \\
\\ 
\texttt{In contrast, Method 2 does not condition on any intermediate variables, so it is less likely to introduce bias due to mediator conditioning.} \\
\\ 
\texttt{Therefore, it is more correct to use Method 2 than Method 1 to understand how maternal smoking status affects freckles.} \\
\\ 
\texttt{References:} \\
\texttt{Pearl, J. (2009). Causality: Models, reasoning, and inference. Cambridge University Press.} \\
            }
        }
    }
    & & \\ \\

    \theutterance \stepcounter{utterance}  
    & & & \multicolumn{2}{p{0.3\linewidth}}{
        \cellcolor[rgb]{0.9,0.9,0.9}{
            \makecell[{{p{\linewidth}}}]{
                \texttt{\tiny{[GM$|$GM]}}
                \texttt{No} \\
            }
        }
    }
    & & \\ \\

    \theutterance \stepcounter{utterance}  
    & & & \multicolumn{2}{p{0.3\linewidth}}{
        \cellcolor[rgb]{0.9,0.9,0.9}{
            \makecell[{{p{\linewidth}}}]{
                \texttt{\tiny{[GM$|$GM]}}
                \texttt{no} \\
            }
        }
    }
    & & \\ \\

    \theutterance \stepcounter{utterance}  
    & & & \multicolumn{2}{p{0.3\linewidth}}{
        \cellcolor[rgb]{0.9,0.9,0.9}{
            \makecell[{{p{\linewidth}}}]{
                \texttt{\tiny{[GM$|$GM]}}
                \texttt{game\_result = WIN} \\
            }
        }
    }
    & & \\ \\

\end{supertabular}
}

\end{document}
